\documentclass[11pt]{article}
\usepackage[margin=1in]{geometry}
\usepackage{setspace}
\usepackage{booktabs}
\usepackage{amsmath,amssymb}
\usepackage{graphicx}
\usepackage{svg}
\usepackage{float}
\usepackage{natbib}
\usepackage[hidelinks]{hyperref}
\graphicspath{{../../results/figures/}}
\svgpath{{../../results/figures/}}
\doublespacing

\title{Household and Community Heterogeneity in Childhood Stunting in Ghana: A Bayesian Hierarchical Analysis of the 2022 Demographic and Health Survey}
\author{}
\date{}

\begin{document}
\maketitle

\begin{abstract}
Childhood stunting remains unevenly distributed in Ghana despite substantial national progress. We analysed 4,928 children aged 0--59 months from the 2022 Ghana Demographic and Health Survey using survey-weighted descriptive methods and Bayesian hierarchical logistic regression with children nested within households and survey communities. Sequential models introduced household and community random intercepts, water and sanitation characteristics, and maternal education. Survey-weighted stunting prevalence was 17.4\% (95\% CI: 15.8--18.9). In the household-and-community model, the posterior median household random-intercept SD was 1.07 compared with 0.42 at the community level. In the final model, boys had higher posterior odds of stunting than girls (OR 1.50, 95\% posterior interval 1.24--1.82), children aged 24--35 months had higher odds than those aged 0--5 months (OR 2.91, 2.06--4.21), unimproved drinking-water source was associated with higher odds (OR 1.50, 1.15--1.99), and higher maternal education was associated with lower odds (OR 0.37, 0.20--0.65). Household heterogeneity remained substantially larger than community heterogeneity after covariate adjustment. Results were robust to alternative priors, missing-data treatment, age functional form, and WASH coding, although water estimates were less precise under survey-weight sensitivity analysis. Residual stunting heterogeneity in Ghana appears concentrated more strongly at the household than community level.
\end{abstract}

\noindent\textbf{Keywords:} child nutrition; stunting; Bayesian hierarchical model; Ghana; Demographic and Health Survey; household factors; water sanitation and hygiene

\section*{Key Messages}
\begin{itemize}
\item Residual childhood-stunting heterogeneity in the 2022 Ghana DHS was substantially larger between households than between survey communities.
\item The household-level variance remained large after adjustment for child characteristics, wealth, region, residence, WASH, and maternal education.
\item Unimproved drinking-water source was associated with higher posterior odds of stunting, with the clearest detailed association observed for surface-water use.
\item Higher maternal education and greater household wealth were associated with lower posterior odds of stunting.
\item The main conclusions were stable across multiple sensitivity analyses, although water estimates became less precise under survey-weighted pseudo-posterior analysis.
\end{itemize}

\section{Introduction}
Childhood stunting reflects chronic or recurrent constraints on linear growth and remains an important marker of nutritional and social disadvantage. Ghana has experienced a substantial long-term decline in stunting, yet the 2022 Ghana Demographic and Health Survey (GDHS) still estimated that approximately 17\% of children under five were stunted \citep{GSSICF2024}. The national average also masks pronounced regional, socioeconomic, and household-level inequalities.

The determinants of stunting operate across multiple levels. Individual characteristics such as age and sex interact with maternal, household, environmental, and broader contextual conditions. Children living in the same household share food resources, caregiving environments, sanitation facilities, socioeconomic shocks, and other exposures, while households located within the same survey community may share health services, infrastructure, markets, and environmental conditions. Analyses that ignore these hierarchical dependencies may underestimate uncertainty and can obscure the level at which residual heterogeneity is concentrated.

Previous Ghanaian studies have established the relevance of multilevel and spatial approaches. Bayesian geostatistical analyses have documented geographical variation in stunting \citep{AhetoDagne2021,AmoakoJohnson2022}, while multilevel studies have identified persistent socioeconomic and household correlates \citep{AhetoEtAl2015,IddrisuGyabaah2023}. Recent work using the 2022 GDHS has also examined WASH and nutritional outcomes among children born to young mothers \citep{OgumEtAl2026}. Bayesian or multilevel modelling is therefore not itself novel in this setting.

The present study focuses instead on a narrower question: how much residual heterogeneity in childhood stunting is concentrated at the household level relative to the community level in the most recent nationally representative GDHS, and how stable is that decomposition after household WASH and maternal education are introduced? We use Bayesian hierarchical logistic regression to separate household and community random effects and evaluate the robustness of substantive conclusions across alternative priors, weighting approaches, missing-data treatment, age functional forms, and WASH coding.

\section{Methods}
\subsection{Data source and analytic sample}
We used the 2022 GDHS Household Member Recode under authorized DHS access. The survey used a stratified two-stage cluster design \citep{GSSICF2024}. The analytic population comprised de facto children aged 0--59 months with valid height-for-age measurements. Children were retained if they had slept in the sampled household the night before interview (\texttt{hv103 = yes}), age was between 0 and 59 months, and the DHS height-for-age z-score variable (\texttt{hc70}) fell within the valid range of -600 to 600.

Stunting was defined as height-for-age z-score below -2 standard deviations, operationalized as \texttt{hc70 < -200}. The final core sample contained 4,928 children nested within 3,545 households and 616 survey communities.

\subsection{Variables}
Core covariates were child sex, age group, household wealth quintile, urban/rural residence, and administrative region. Age was grouped as 0--5, 6--11, 12--23, 24--35, 36--47, and 48--59 months in the primary analysis.

Household drinking-water source (\texttt{hv201}) and sanitation facility (\texttt{hv205}) were classified as improved versus unimproved/no facility according to JMP-consistent source/facility-type definitions \citep{WHOUNICEFJMP2021}. These categories represent facility type rather than complete JMP service levels. Maternal educational attainment was obtained from \texttt{hc61} and classified as no education, primary, secondary, or higher. Maternal education was missing for 425 children, yielding a 4,503-child complete-case sample for the final maternal-education model.

\subsection{Survey-weighted descriptive analysis}
Descriptive prevalence estimates used household-member sampling weights (\texttt{hv005}). Taylor-linearized standard errors and 95\% confidence intervals accounted for primary sampling units (\texttt{hv021}) and sampling strata (\texttt{hv022}). The weighted analytic denominator was 4,292.97, matching the official GDHS weighted height-for-age denominator of 4,293 after rounding.

\subsection{Bayesian hierarchical models}
For child $i$ in household $j$ and community $k$,
\[
Y_{ijk}\sim\text{Bernoulli}(p_{ijk}),
\]
with
\[
\text{logit}(p_{ijk})=\alpha+\mathbf{x}_{ijk}^{\top}\boldsymbol{\beta}+u_k+v_{jk},
\]
where $u_k\sim N(0,\tau_c^2)$ is a community random intercept and $v_{jk}\sim N(0,\tau_h^2)$ is a household random intercept. Non-centred parameterizations were used for both random effects.

The main priors were $\alpha\sim N(0,1.5^2)$, $\beta_r\sim N(0,1^2)$, and $\tau_c,\tau_h\sim\text{Half-Normal}(1)$. Model 1 contained the core covariates and both household and community random intercepts. Model 2 added WASH indicators. Model 3 added maternal education. To avoid conflating a covariate effect with a sample change, Model 2 was also refitted on the same 4,503-child complete-maternal-information sample used by Model 3.

\subsection{Hierarchical variance measures}
For the logistic model, the latent residual variance was taken as $\pi^2/3$. The community ICC and household-specific variance partition coefficient were
\[
ICC_c=\frac{\tau_c^2}{\tau_c^2+\tau_h^2+\pi^2/3},
\qquad
VPC_h=\frac{\tau_h^2}{\tau_c^2+\tau_h^2+\pi^2/3}.
\]
Median odds ratios were calculated as
\[
MOR=\exp\left[\sqrt{2}\Phi^{-1}(0.75)\tau\right].
\]

\subsection{Diagnostics and robustness analyses}
The principal models were fitted with NUTS in PyMC. The strengthened final Model 3 diagnostic run combined eight independent chains with 500 warmup and 500 retained draws per chain. We examined divergences, $\hat R$, effective sample size, BFMI, and tree-depth behavior. Posterior predictive checks compared replicated with observed overall and age-specific prevalence.

Sensitivity analyses evaluated tighter and wider prior scales, a survey-weight pseudo-posterior with weights normalized to mean one, explicit inclusion of a missing maternal-education category, a cubic B-spline representation of child age, and more detailed WASH categories. Predictive comparison of harmonized Models 2 and 3 used PSIS-LOO with Pareto-$k$ diagnostics.

\subsection{Ethics and data access}
This study used de-identified secondary data obtained through authorized access from The DHS Program. The original GDHS obtained informed consent under its approved survey procedures. Restricted DHS microdata are not redistributed with the manuscript or public repository. Public reproducibility materials contain analysis code and aggregate, non-identifying outputs only.

\section{Results}
\subsection{Descriptive results}
The analytic sample included 927 stunted children among 4,928 children. Survey-weighted stunting prevalence was 17.4\% (95\% CI: 15.8--18.9). Prevalence was 19.4\% among boys and 15.3\% among girls. It peaked at 24.6\% among children aged 24--35 months. Prevalence declined from 22.8\% in the poorest wealth quintile to 7.5\% in the richest. Regionally, the highest weighted prevalence estimates were observed in Northern (29.6\%) and North East (29.3\%) regions.

\begin{figure}[H]
\centering
\includesvg[width=0.90\textwidth]{stunting_by_region}
\caption{Survey-weighted prevalence of childhood stunting by region.}
\end{figure}

\subsection{Household and community heterogeneity}
Model 1 revealed substantially greater residual heterogeneity between households than between communities. The posterior median community SD was 0.416 (95\% posterior interval 0.159--0.590), whereas the household SD was 1.066 (0.774--1.346). The median community ICC was 0.038 and household VPC was 0.246. The corresponding MORs were 1.49 for communities and 2.76 for households.

This contrast persisted after WASH and maternal education were introduced. In the strengthened eight-chain Model 3 run, the posterior median community SD was 0.389 and the household SD was 1.234. The community ICC was 0.030 and household VPC 0.307, with MORs of approximately 1.45 and 3.24, respectively.

\subsection{Fixed effects}
In the final model, boys had higher posterior odds of stunting than girls (OR 1.50, 95\% posterior interval 1.24--1.82). Children aged 24--35 months had substantially higher odds than those aged 0--5 months (OR 2.91, 2.06--4.21). The richest-versus-poorest posterior OR was 0.36 (0.20--0.64).

Unimproved drinking-water source was associated with higher posterior odds of stunting (OR 1.50, 1.15--1.99). The sanitation estimate was smaller and more uncertain (OR 1.12, 0.86--1.46). Higher maternal education was associated with substantially lower posterior odds of stunting compared with no formal education (OR 0.37, 0.20--0.65). Primary (OR 1.02, 0.75--1.41) and secondary education (OR 0.89, 0.68--1.14) were not clearly separated from no education after full adjustment.

\begin{figure}[H]
\centering
\includesvg[width=0.94\textwidth]{model3_posterior_or_forest}
\caption{Selected posterior odds ratios from the final maternal-education model.}
\end{figure}

\subsection{Diagnostics and robustness}
The strengthened Model 3 run had zero divergences, BFMI values between 0.51 and 0.66, and no tree-depth saturation. Fixed-effect convergence was excellent, with maximum $\hat R$ around 1.003 and bulk effective sample sizes above 2,200. The community and household SDs remained slower-mixing, with $\hat R$ approximately 1.021 and 1.018 and bulk ESS approximately 544 and 524, respectively.

Posterior predictive checks reproduced the observed overall prevalence and all six age-specific prevalences. The age-spline sensitivity produced nearly identical substantive estimates. Under more detailed WASH coding, surface-water use showed the clearest elevated association relative to improved water (OR 1.52, 1.10--2.08). The survey-weight pseudo-posterior preserved the direction of the main results but widened uncertainty for water: OR 1.38 (0.90--2.07).

PSIS-LOO showed essentially equivalent predictive performance between harmonized Model 2 and Model 3. Model 3 improved ELPD by only 0.92 (SE 3.38). Model 3 had no observations with Pareto $k>0.70$; Model 2 had five, and neither model had Pareto $k>1$.

\section{Discussion}
The central finding is that residual heterogeneity in childhood stunting is much larger between households than between survey communities. This pattern persists after adjustment for child characteristics, household wealth, residence, region, WASH, and maternal education. Earlier Ghanaian research has also identified meaningful household-level variation in childhood malnutrition \citep{AhetoEtAl2015}, while geostatistical analyses have documented broader spatial heterogeneity \citep{AhetoDagne2021,AmoakoJohnson2022}. The present results show that these two levels should not be conflated.

The magnitude of the household MOR illustrates the substantive importance of this heterogeneity. In the final model, the household MOR exceeded three, whereas the community MOR was approximately 1.45. This suggests that otherwise similar children can experience substantially different latent risk depending on their household context. Such variation may reflect food availability and allocation, caregiving practices, maternal health, household infection environments, economic shocks, or other unmeasured shared exposures.

The age pattern is consistent with accumulated growth faltering during the transition from infancy through the second and third years of life. The robustness of the fixed effects under a spline representation of age indicates that the broader conclusions are not driven by the chosen age categories. The elevated odds among boys are also consistent with broader sub-Saharan African evidence \citep{TakeleEtAl2022}.

The persistent wealth gradient is consistent with long-standing socioeconomic inequality in child nutrition in Ghana \citep{SeworJayalakshmi2024,OsborneEtAl2025}. Its attenuation after maternal education was introduced is consistent with overlap between education and household socioeconomic position, although the cross-sectional design does not support causal mediation claims.

The WASH results require a more nuanced interpretation than a simple improved/unimproved contrast. The detailed sensitivity analysis suggests that the water association is driven most clearly by surface-water exposure. In contrast, sanitation associations remained uncertain after socioeconomic and hierarchical adjustment. The primary water association also became less precise under the survey-weight pseudo-posterior.

Higher maternal education showed a strong inverse association with stunting. Similar associations have been reported in Ghana and across sub-Saharan Africa \citep{AhetoEtAl2015,TakeleEtAl2022,AgyenEtAl2024}. The lack of clear independent associations for primary and secondary education after full adjustment argues against a simplistic dose-response interpretation.

The predictive comparison is also informative. Model 3 adds an interpretable maternal-education association but offers little improvement in out-of-sample predictive performance relative to harmonized Model 2. This distinction between explanatory interpretation and prediction should be preserved rather than overstating the value of the larger model.

\subsection{Strengths and limitations}
Strengths include use of the most recent nationally representative GDHS, exact reproduction of the official weighted stunting estimate, explicit separation of household and community random effects, same-sample model comparisons, and extensive robustness checks.

Several limitations remain. The cross-sectional design precludes causal inference. Some relevant determinants, including dietary intake, food insecurity, maternal nutritional status, birth size, and repeated infection history, were not included. The WASH variables describe source/facility type rather than full JMP service levels. Maternal education was missing for 8.6\% of the core sample, although a full-sample sensitivity analysis yielded similar conclusions.

The primary Bayesian likelihood does not fully encode the complex survey design; survey-weighted pseudo-posterior results are therefore used as a sensitivity analysis rather than a replacement for the main hierarchical model. Many households contributed only one eligible child, so household variance remains a model-based population parameter rather than a directly observed correlation for every household.

Finally, the hierarchical SDs continue to mix more slowly than the fixed effects even in the strengthened eight-chain run. Their effective sample sizes are above 500 and the broad household-versus-community contrast is stable, but exact variance-component intervals should still be interpreted with more caution than the fixed-effect estimates.

\section{Conclusion}
Childhood stunting in Ghana remains strongly patterned by age, sex, socioeconomic position, household environmental conditions, and maternal education. The most distinctive finding is that unexplained heterogeneity is concentrated substantially more strongly at the household than the community level. This pattern persists after WASH and maternal education are introduced.

The findings support closer attention to shared household conditions in child-nutrition research and programming while avoiding causal interpretations that the cross-sectional data cannot sustain.

\section*{Data Availability Statement}
The 2022 Ghana DHS microdata are owned and distributed by The DHS Program and cannot be redistributed by the authors. Eligible researchers may request access directly from The DHS Program subject to its data-use conditions. Reproducibility code, aggregate tables, figures, and non-identifying diagnostic outputs are available in the project repository; restricted row-level DHS data are not included.

\section*{Conflict of Interest Statement}
The authors declare no conflict of interest, subject to confirmation by all co-authors before submission.

\section*{Funding}
[INSERT FUNDING STATEMENT AFTER AUTHOR CONFIRMATION.]

\bibliographystyle{apalike}
\bibliography{../../thesis/references}
\end{document}
